\documentclass[aspectratio=169,11pt]{beamer}

\usetheme{Madrid}
\usecolortheme{default}
\usefonttheme{professionalfonts}
\setbeamertemplate{navigation symbols}{}
\setbeamertemplate{footline}[frame number]

\usepackage{amsmath, amssymb, booktabs, graphicx}

\newcommand{\E}{\mathbb{E}}
\newcommand{\Cov}{\operatorname{Cov}}
\newcommand{\1}{\mathbf{1}}

\title{Instrumental Variables, 2SLS, and Instrument Design}
\subtitle{Empirical Methods --- Week 4}
\author{Teaching deck draft}
\date{}

\begin{document}

\begin{frame}
  \titlepage
\end{frame}

\begin{frame}{Week 4 question}
\begin{itemize}
    \item When can an external source of variation identify a causal effect?
    \item What treatment margin does the instrument move?
    \item Why should the instrument affect outcomes only through treatment?
    \item For whom is the resulting estimate interpretable?
\end{itemize}

\medskip
IV is a design argument before it is a two-stage regression.
\end{frame}

\begin{frame}{Why IV is needed}
\[
Y_i = \beta D_i + X_i'\gamma + u_i
\]

\begin{itemize}
    \item OLS fails when $D_i$ is correlated with unobserved determinants of $Y_i$.
    \item Endogeneity can come from ability, sorting, reverse causality, anticipation, or measurement error.
    \item IV asks for a source of variation $Z_i$ that shifts treatment but is otherwise excluded from outcomes.
    \item The instrument's economic content defines the design.
\end{itemize}
\end{frame}

\begin{frame}{First stage, reduced form, Wald}
\[
D_i = \pi Z_i + X_i'\delta + v_i
\qquad
Y_i = \rho Z_i + X_i'\kappa + \varepsilon_i
\]

\[
\beta^{IV}
=
\frac{\E[Y_i\mid Z_i=1]-\E[Y_i\mid Z_i=0]}
{\E[D_i\mid Z_i=1]-\E[D_i\mid Z_i=0]}
=
\frac{\Cov(Y_i,Z_i)}{\Cov(D_i,Z_i)}
\]

\begin{itemize}
    \item Numerator: instrument effect on outcome.
    \item Denominator: instrument effect on treatment.
    \item Ratio: effect of the treatment margin moved by the instrument.
\end{itemize}
\end{frame}

\begin{frame}{2SLS projection logic}
\[
W=[X,Z],
\qquad
P_W=W(W'W)^{-1}W',
\qquad
\widehat{D}=P_WD
\]

\[
\widehat{\theta}_{2SLS}
=
(\widetilde X'P_W\widetilde X)^{-1}\widetilde X'P_WY,
\qquad
\widetilde X=[D,X]
\]

\begin{itemize}
    \item First stage keeps the treatment variation predicted by instruments and controls.
    \item Second stage relates outcomes to that instrument-predicted treatment.
    \item Projection notation does not create identification; the instrument design does.
\end{itemize}
\end{frame}

\begin{frame}{Assumptions}
\small
\begin{tabular}{p{0.20\linewidth} p{0.35\linewidth} p{0.34\linewidth}}
\toprule
Assumption & Meaning & Main design question \\
\midrule
Relevance & $Z$ shifts $D$ & What treatment margin moves? \\
Independence & $Z$ is as-good-as-random given $X$ & What assignment or shock process creates variation? \\
Exclusion & $Z$ affects $Y$ only through $D$ & What direct channels remain? \\
Monotonicity & no defiers & Could $Z$ push some units away from treatment? \\
First-stage content & compliers are interpretable & Are they the margin the paper cares about? \\
\bottomrule
\end{tabular}
\end{frame}

\begin{frame}{LATE and compliers}
\[
\beta^{IV}
=
\E\left[Y_i(1)-Y_i(0)\mid D_i(1)>D_i(0)\right]
\]

\begin{itemize}
    \item Compliers take treatment when exposed to the instrument and not otherwise.
    \item Always-takers and never-takers shape take-up levels, but not the Wald numerator-denominator margin.
    \item The estimate can be policy-relevant when compliers are the policy margin.
    \item It should not be reported as a universal ATE unless the design justifies that move.
\end{itemize}
\end{frame}

\begin{frame}{Weak instruments and many instruments}
\[
F_{\text{excluded }Z}
=
\frac{(RSS_R-RSS_U)/q}{RSS_U/(n-k_U)}
\]

\begin{itemize}
    \item Weak first stages make 2SLS noisy, biased toward OLS, and specification-sensitive.
    \item Report first-stage coefficients, economic magnitude, and weak-IV diagnostics.
    \item Use weak-IV robust inference when the first stage is fragile.
    \item Many instruments can overfit the first stage and make interpretation opaque.
    \item Multiple instruments may combine different complier populations.
\end{itemize}
\end{frame}

\begin{frame}{Instrument design gallery}
\small
\begin{tabular}{p{0.24\linewidth} p{0.34\linewidth} p{0.31\linewidth}}
\toprule
Family & First-stage logic & Core threat \\
\midrule
Policy rules & eligibility or legal exposure shifts treatment & cohort, seasonality, or rule channels \\
Distance/access & lower cost raises take-up & geography directly affects outcomes \\
Leniency leave-out & assigned decision-makers differ in propensities & decision-makers affect outcomes through other channels \\
Shift-share & shocks hit exposed units differently & endogenous shares or non-random shocks \\
Historical legacy & past conditions predict current treatment & many persistent pathways to current outcomes \\
\bottomrule
\end{tabular}

\medskip
Every family needs its own first-stage story, exclusion argument, and inference level.
\end{frame}

\begin{frame}{Shift-share instruments}
\[
Z_{\ell}=\sum_s w_{\ell s}g_s
\]

\begin{itemize}
    \item Shares $w_{\ell s}$ measure exposure; shocks $g_s$ create differential predicted treatment.
    \item Modern practice asks whether identification comes from quasi-random shocks or as-good-as-random shares.
    \item Endogenous shares can proxy local trends, institutions, or industry composition.
    \item Inference should reflect the shock structure, not only geographic clustering.
    \item Report dominant shocks, leave-out choices, exposure balance, and shock-level sensitivity.
\end{itemize}
\end{frame}

\begin{frame}{Leniency leave-out instruments}
\[
Z_{ij}
=
\frac{1}{N_j-1}\sum_{k\neq i:J_k=j}D_k
\]

\begin{itemize}
    \item The leave-out measure avoids using case $i$ to predict its own treatment.
    \item Identification requires quasi-random assignment to decision-makers within cells.
    \item Balance checks should be done inside the assignment cells that define randomization.
    \item Exclusion can fail if lenient decision-makers change information, processing time, stigma, or other treatments.
    \item The effect is local to marginal cases moved by decision-maker style.
\end{itemize}
\end{frame}

\begin{frame}{Historical, distance, and policy instruments}
\begin{itemize}
    \item Policy rules: compelling when the rule is sharp, binding, and isolated from other channels.
    \item Distance/access: useful when access cost is the treatment margin; fragile when geography carries opportunity or family background.
    \item Historical instruments: powerful first stages, but persistent direct channels are the central threat.
    \item The implementation task is to name the mechanism, show first-stage content, and stress-test exclusion.
\end{itemize}
\end{frame}

\begin{frame}{Research Lab design}
\begin{columns}[T,onlytextwidth]
\begin{column}{0.32\linewidth}
\textbf{Reproduce}
\begin{itemize}
    \item Compulsory-schooling synthetic IV
    \item First stage
    \item Reduced form
    \item Wald and 2SLS
\end{itemize}
\end{column}
\begin{column}{0.32\linewidth}
\textbf{Diagnose}
\begin{itemize}
    \item Partial first-stage $F$
    \item Balance by instrument
    \item Complier profile
    \item Exclusion memo
\end{itemize}
\end{column}
\begin{column}{0.32\linewidth}
\textbf{Transfer}
\begin{itemize}
    \item Shift-share construction
    \item Shock-level diagnostics
    \item Endogenous shares
    \item Leniency comparison
\end{itemize}
\end{column}
\end{columns}
\end{frame}

\begin{frame}{Takeaways}
\begin{itemize}
    \item IV estimates the treatment margin induced by the instrument.
    \item The first stage defines the compliers and the economic interpretation.
    \item Exclusion is the central substantive assumption and is not directly testable.
    \item Weak and many instruments can make formal 2SLS output misleading.
    \item Instrument families are not interchangeable; each has distinct implementation caveats.
    \item A credible IV paper is a theory of variation before it is a regression table.
\end{itemize}
\end{frame}

\end{document}
